\documentclass[a4paper,12pt]{article}
\usepackage[T2A]{fontenc}
\usepackage[utf8]{inputenc}
\usepackage{cmap}
\usepackage[english, russian]{babel}
\usepackage{wrapfig}
\usepackage[backend=biber,style=gost-numeric]{biblatex} % стиль ГОСТ
\addbibresource{literatura.bib}
\usepackage{misccorr} % в заголовках появляется точка, но при ссылке на них ее нет
\usepackage{amssymb,amsfonts,amsmath,amsthm}  
\usepackage{indentfirst}
\usepackage[usenames,dvipsnames]{color} 
\usepackage[unicode,hidelinks]{hyperref}
% \hypersetup{%
%     pdfborder = {0 0 0}
\usepackage{multirow}
\usepackage{makecell,multirow} 
\usepackage{ulem}
%\usepackage{hyperref}
%\usepackage{graphicx,wrapfig}
%\renewcommand\epsilon{\varepsilon}
% \renewcommand\epsilon{\varepsilon}
%\graphicspath{{img/}}
\usepackage{geometry}
\geometry{left=1cm,right=1cm,top=2cm,bottom=2cm,bindingoffset=0cm,headheight=15pt}
\usepackage{fancyhdr} 
\linespread{1.2} 
\frenchspacing 
\renewcommand{\labelenumii}{\theenumii)} 
% \usepackage{caption}
%%%%%%%%%%%%%%%%%%%%%%%%%%%%%%%%%%%%%%%%%%%%%%%%%%%%%%%%%%%%%%%%%%%%%%%%%%%%%%%
%%%%%%%%%%%%%%%%%%%%%%%%%%%%%%%%%%%%%%%%%%%%%%%%%%%%%%%%%%%%%%%%%%%%%%%%%%%%%%%

\def\labauthor{Сарафанов И.Г.}
\def\labauthors{\labauthor}
\def\labnumber{127}
\def\labtheme{Таблица производных и интегралов}

%%%%%%%%%%%%%%%%%%%%%%%%%%%%%%%%%%%%%%%%%%%%%%%%%%%%%%%%%%%%%%%%%%%%%%%%%%%%%%%
	%применим колонтитул к стилю страницы
\pagestyle{fancy} 
	%очистим "шапку" страницы
\fancyhead{\LaTeX} 
	%слева сверху на четных и справа на нечетных
\fancyhead[L]{\labauthors} 
	%справа сверху на четных и слева на нечетных
\fancyhead[R]{Иродов \textnumero1.137} 
	%очистим "подвал" страницы
\fancyfoot{} 
	% номер страницы в нижнем колинтуле в центре
\fancyfoot[C]{\thepage} 
%%%%%%%%%%%%%%%%%%%%%%%%%%%%%%%%%%%%%%%%%%%%%%%%%%%%%%%%%%%%%%%%%%%%%%%%%%%%%%%
\usepackage{float}
\usepackage[mode=buildnew]{standalone}
\usepackage{tikz} 
% \usepackage{subcaption}
\usepackage{tikz,csvsimple,physics}
\makeatother
%\usepackage{booktabs}
\usepackage{pgfplots, pgfplotstable}
\usepackage[outline]{contour}
\usepackage{tocloft}
\renewcommand{\cftsecleader}{\cftdotfill{\cftdotsep}}


% вот этот удивительный мир

\begin{document}
% Задание 3389
\section*{3389}
Найти \(\displaystyle \frac{\partial^2 z}{\partial x^2},\; \frac{\partial^2 z}{\partial x \partial y},\; \frac{\partial^2 z}{\partial y^2}\) при \(x = 1,\; y = -2,\; z = 1\), если
\[ x^2 + 2y^2 + 3z^2 + xy - z - 9 = 0. \]

% Задание 3410
\section*{3410}
Пусть функция \(z = z(x, y)\) определяется системой уравнений:
\[
\begin{cases} 
x = e^{u + v}, \\ 
y = e^{u - v}, \\ 
z = uv 
\end{cases}
\]
(\(u\) и \(v\) — параметры). Найти \(dz\) и \(d^2z\) при \(u = 0\) и \(v = 0\).

% Задание 3411
\section*{3411}
Найти \(\displaystyle \frac{\partial z}{\partial x}\) и \(\displaystyle \frac{\partial^2 z}{\partial x^2}\), если
\[z = x^2 + y^2,\]
где \(y = y(x)\) определяется из уравнения
\[x^2 - xy + y^2 = 1.\]

% Задание 3412
\section*{3412}
Найти \(\displaystyle \frac{\partial u}{\partial x}\) и \(\displaystyle \frac{\partial u}{\partial y}\), если
\[u = \frac{x + z}{y + z},\]
где \(z\) определяется из уравнения
\[ze^z = xe^x + ye^y.\]

% Задание 3413
\section*{3413}
Пусть уравнения
\[x = \varphi(u, v),\; y = \psi(u, v),\; z = \chi(u, v)\]
определяют \(z\) как функцию от \(x\) и \(y\). Найти \(\displaystyle \frac{\partial z}{\partial x}\) и \(\displaystyle \frac{\partial z}{\partial y}\).

% Задание 3414
\section*{3414}
Пусть
\[x = \varphi(u, v),\; y = \psi(u, v).\]
Найти частные производные первого и второго порядков от обратных функций: \(u = u(x, y)\) и \(v = v(x, y)\).

% Задание 3415
\section*{3415}
Найти \(\displaystyle \frac{\partial u}{\partial x},\; \frac{\partial u}{\partial y},\; \frac{\partial v}{\partial x},\; \frac{\partial v}{\partial y}\), если:
\begin{enumerate}
    \item[a)] \(x = u \cos \dfrac{v}{u},\; y = u \sin \dfrac{v}{u};\)
    \item[б)] \(x = e^u + u \sin v,\; y = e^u - u \cos v.\)
\end{enumerate}
\end{document}